\SourceLeaf{295}%
and the mean of the two azimuths. If we write
\[
\tfrac{1}{2}(S+S')=B,\quad\tfrac{1}{2}(T+T')=A,\quad S-S'=b,\quad T-T'=a
\]
we first obtain the formulas
\[
\begin{aligned}
\sin\tfrac{1}{2}r\sin A&=\sin\tfrac{1}{2}\lambda\cos B\\
\sin\tfrac{1}{2}r\cos A&=\cos\tfrac{1}{2}\lambda\sin\tfrac{1}{2}b\\
\cos\tfrac{1}{2}r\sin\tfrac{1}{2}a&=\sin\tfrac{1}{2}\lambda\sin B\\
\cos\tfrac{1}{2}r\cos\tfrac{1}{2}a&=\cos\tfrac{1}{2}\lambda\cos\tfrac{1}{2}b
\end{aligned}
\]
Thus, if \(A,B,r\) are taken as given, we determine \(a\) and \(\lambda\) from the formulas
\[
\begin{aligned}
\sin A\operatorname{tang}B\operatorname{tang}\tfrac{1}{2}r&=\sin\tfrac{1}{2}a\\
\frac{\sin A\sin\tfrac{1}{2}r}{\cos B}&=\sin\tfrac{1}{2}\lambda
\end{aligned}
\]
and then \(b\) from
\[
\frac{\cos A\operatorname{tang}\tfrac{1}{2}r}{\cos\tfrac{1}{2}a}=\operatorname{tang}\tfrac{1}{2}b
\]
or from
\[
\frac{\cos A\sin\tfrac{1}{2}r}{\cos\tfrac{1}{2}\lambda}=\sin\tfrac{1}{2}b
\]
Because \(r,a,\lambda,b\) are small, however, we will prefer to use the following formulas. They are much more convenient and are accurate up to but excluding fifth order:
\[
\begin{aligned}
a^0&=r\sin A\operatorname{tang}B\\
\lambda^0&=\frac{r\sin A}{\cos B}\\
b^0&=r\cos A\\
\log a&=\log a^0+\mu rr+\tfrac{1}{2}\mu a^0a^0\\
\log\lambda&=\log\lambda^0-\tfrac{1}{2}\mu rr+\tfrac{1}{2}\mu\lambda^0\lambda^0\\
\log b&=\log b^0+\tfrac{1}{2}\mu a^0a^0+\mu\lambda^0\lambda^0
\end{aligned}
\]
As we can see, the third correction equals the sum of the first and twice the second.

These formulas give no direct solution to our problem, but they can serve as a check or as a concise overview of the direct solution. Anyone with some skill in numerical calculations will also find it easy to use them for an indirect solution and will prefer that solution to all others for its convenience and accuracy, especially when a rough approximation has already been made for other purposes.
\par\medskip\centerline{\rule{.25\textwidth}{.3pt}}
